\chapter{Endovenous Ablation}

\section*{Overview}
Endovenous ablation closes a superficial leg vein whose valves allow blood to flow backward (reflux), often contributing to varicose veins. A clinician guides a thin catheter into the vein using ultrasound, then closes it from the inside. Methods include heat from laser or radiofrequency energy and nonthermal methods such as medical adhesive. Blood is redirected through other healthy veins. The treated vein gradually becomes a closed cord of scar tissue; it is not usually removed.

\section*{Alternative Names}
Endovenous thermal ablation; endovenous laser ablation (EVLA or EVLT); radiofrequency ablation (RFA); endovenous vein closure. Nonthermal options include cyanoacrylate adhesive closure and mechanochemical ablation. VenaSeal is a brand name for one adhesive system.

\section*{Principle}
With thermal ablation, heat injures the vein wall so the vein closes. Local anesthetic fluid is placed around the vein to numb it and protect nearby tissue. With adhesive closure, a medical glue seals the vein; mechanochemical methods combine mechanical action with a chemical agent. The specific steps, benefits, and risks differ by method.

\section*{History}
Endovenous thermal techniques became widely used from the late 1990s onward as less invasive alternatives to surgical vein stripping for many people with reflux in the main superficial leg veins. Nonthermal catheter methods are also available. Choice depends on vein anatomy, symptoms, patient factors, and clinician experience.

\section*{Why It Is Done}
The procedure may be offered for confirmed superficial venous reflux that causes troublesome varicose veins, aching, heaviness, swelling, itching, skin changes, or venous ulcers. It is intended to treat the faulty source vein and relieve symptoms or improve venous disease. It does not prevent new varicose veins from developing in other veins.

\section*{Role in Care}
It is a treatment, not a diagnostic test. A duplex ultrasound scan is used to confirm reflux, map the veins, and plan treatment. For suitable people with symptomatic varicose veins and truncal reflux, guidelines commonly recommend endothermal ablation as an initial treatment option. Foam sclerotherapy, surgery, adhesive closure, or other options may be considered depending on anatomy, suitability, availability, and preference.

\section*{How It Is Performed}
The procedure is usually done in an outpatient setting with ultrasound guidance. The clinician numbs the skin and accesses the vein through a small puncture. A catheter is advanced to the planned location. For laser or radiofrequency treatment, anesthetic fluid is placed around the vein and energy is delivered as the catheter is withdrawn. For adhesive or other nonthermal methods, the vein is closed using the method-specific device, and the amount of local anesthetic may differ. Additional procedures, such as removing bulging side branches through tiny incisions or treating them with injections, may be done at the same visit or later.

\section*{Preparation}
The clinic explains the chosen method, medicines, and aftercare. Duplex ultrasound maps the veins before treatment. Tell the team about previous blood clots, medicines (especially blood thinners), allergies, pregnancy or possible pregnancy, and any reduced blood flow to the legs. Follow instructions about eating, drinking, and medicines; these depend on the method and whether sedation is planned. Do not stop prescribed medicines unless the treating team advises you to.

\section*{Contraindications and Precautions}
An active deep vein thrombosis or infection in the treatment area generally needs assessment before an elective ablation. A history of DVT does not automatically rule out treatment, but may affect planning and clot-risk precautions. Deep venous obstruction, pregnancy, limited mobility, significant arterial disease, and the exact vein anatomy may change whether, when, or how treatment is offered. The clinician reviews these factors and discusses alternatives. Compression advice also needs to account for arterial circulation.

\section*{Risks}
Common temporary effects include bruising, tenderness, and a firm or sensitive area along the treated vein. Other risks include skin irritation or burns with thermal methods, numbness or tingling from nearby nerve irritation, inflammation of a superficial vein, and skin color changes. A blood clot can form in a deep vein or extend toward it; this is uncommon but potentially serious. Rare complications include infection or injury to nearby structures. Varicose veins or symptoms can persist or return, and additional treatment may be needed.

\section*{Aftercare and Recovery}
Most people walk soon after treatment and return to usual light activity quickly. The clinic gives advice about walking, exercise, bathing, compression stockings or bandages, and pain relief. Compression recommendations vary by method and local guidance; it is not a universal long-term requirement after every ablation. Follow the individual instructions you receive. Bruising, tenderness, or a pulling sensation can last for days or weeks as the vein closes.

\section*{Outcome and Follow-up}
The treated vein is expected to close, which can improve symptoms and reduce the pressure feeding varicose veins. Symptom improvement and cosmetic change vary, and visible side branches may need separate treatment. A duplex scan is used when symptoms, clot risk, the type of treatment, or the clinician's plan make it appropriate; early ultrasound is not routinely needed for every average-risk person who feels well after thermal ablation. The care team will explain whether and when you need a scan.

\section*{Expected Results}
The goal is to close the refluxing vein and improve symptoms or complications. The vein may remain palpable for a time as it scars down, and varicose veins may not disappear immediately. Treatment does not cure the tendency to develop venous reflux, so new veins or recurrent symptoms can occur later.

\section*{Follow-up}
Attend follow-up as advised, and report symptoms that persist or return. Further treatment may include ultrasound-guided foam sclerotherapy, phlebectomy, or another vein procedure. Seek assessment for new one-sided leg swelling or pain, since a blood clot needs to be ruled out.

\section*{When to Seek Medical Care}
Contact the treating clinic for increasing redness, warmth, pain, or swelling around the treated area, or for concerns about the puncture site. Seek urgent medical care for sudden one-sided leg swelling or pain, coughing blood, chest pain, fainting, or shortness of breath, which can be signs of a serious blood clot. Follow any additional warning signs in your discharge instructions.

\section*{References}
\begin{enumerate}
  \item National Institute for Health and Care Excellence. Varicose veins: diagnosis and management. https://www.nice.org.uk/guidance/cg168/chapter/Recommendations
  \item Society for Vascular Surgery, American Venous Forum, and American Vein and Lymphatic Society. Clinical Practice Guidelines for the Management of Varicose Veins of the Lower Extremities, Part II. https://pmc.ncbi.nlm.nih.gov/articles/PMC11523430/
  \item Guy's and St Thomas' NHS Foundation Trust. Varicose veins: having a non-surgical procedure. https://www.guysandstthomas.nhs.uk/health-information/varicose-veins/having-non-surgical-procedure-varicose-veins
  \item Healthdirect Australia. Endovenous ablation. https://www.healthdirect.gov.au/procedures/surgery/endovenous-ablation
\end{enumerate}

\clearpage
